


\section{Additional Related Works}

Now we review some additional related works.


\paragraph{PIR with linear computation cost.}
The classical single-server PIR setting does not have preprocessing and the server is required to store the database without further encodings.
Beimel, Ishai and Malkin~\cite{bim-pir} proved that the per-query time of any PIR scheme in the classical model has to be $\Omega(n)$. 
%Inspired by their work, many subsequent studies followed the ``pre-processing'' model to achieve amortized sublinear per-query time.
Early theoretical works~\cite{pir05,pir02,pir06} improve the communication cost of PIR to polylogarithmic. 
Some later works~\cite{xpir,fastpir,spiral} focus on improving the concrete computation time, mainly by reducing the number of public key cryptographic operations. 
%Spiral~\cite{spiral} achieves the best performance 

Within the scope of linear-computation cost PIR, recent schemes~\cite{psir,onionpir,frodopir,simplepir} also utilize preprocessing to improve the concrete performance.
In particular, SimplePIR~\cite{simplepir} achieves the fastest computation time and claims that the online query time in their implementation is already limited by the server's memory I/O speed, instead of the computation.

Nonetheless, all these schemes have linear online time per query, which hinders them from being practical when scaling to large databases. 

\paragraph{Batch PIR.} 
%l~, Batch PIR schemes~\cite{batch-pir-osdi16, angel-pir,vector-pir, pirana} are designed for batched queries.
In batch PIR schemes~\cite{batchcode,batch-pir-osdi16,angel-pir,vector-pir,pirana}, the client is assumed to have a $Q$-size batch of queries simultaneously. 
The client can make a batched query and the server's computation cost for the whole batch is still $O(n)$, while the amortized computation cost for each query is reduced to $\widetilde{O}(n/Q)$. 
The main drawback of batch PIR schemes is that the client cannot make adaptive queries. 
Also, batch PIR schemes usually need some form of homomorphic encryption evaluation on the whole database, making the query latency significantly large.
See the comparison in~\cite{simplepir}.

\paragraph{Preprocessing PIR with sublinear computation cost.} 
Biemel, Ishai and Malkin~\cite{bim-pir} introduced PIR with preprocessing in order to circumvent the $\Omega(n)$ lower-bound on server side computation. They showed that the servers could compute polynomially many bits offline that reduces the online computation for the servers. Specifically, for any $\epsilon>0,k\ge 2$, they gave a $k$-server protocol that has $O(n^{1/2k-1})$ communication and $O(n/(\epsilon\log n)^{2k-2})$ online server computation, and needs $O(n^{1+\epsilon})$ extra server storage. In order to reduce the online server computation further, they gave a second construction which for any $\epsilon>0$ achieves $O(n^{1/2+\epsilon})$ online server computation and communication, and requires $O(n^{1+\epsilon'})$ extra server storage, where $\epsilon'$ depends on $\epsilon$. While the second construction reduces online server computation significantly, the extra server storage is concretely very large, e.g., to achieve $n^{0.6}$ server side computation, the server must store roughly $n^{3.2}$ extra bits. In practice, this kind of overhead is prohibitively large.

Our scheme can be categorized as the ``client-preprocessing'' model, where the client will store some database-dependent hints and use those hints to make efficient online query. 
In fact, most of the recent schemes follow the same design framework proposed by Corrigan-Gibbs and Kogan~\cite{sublinearpir}: the client will store many sets and their corresponding parities as hints and later on the client will find a local set $S$ that contains the online query $x$, then make a carefully-designed query to learn the parity of $S/\{x\}$. 
The initial scheme~\cite{sublinearpir} only supports one query, and subsequent works extend the idea to support multiple queries. 
In the two-server setting, Shi et al.~\cite{shi-pir}, Checklist~\cite{checklist} and TreePIR~\cite{treepir} improve the communication cost to $\widetilde{O}_\lambda(1)$ while the amortized computation cost is $\widetilde{O}_\lambda(\sqrt{n})$. TreePIR is also concretely efficient.
In the single-server setting, Corrigan-Gibbs, Henzinger and Kogan~\cite{pir-ec22} provide a scheme that has $\widetilde{O}_\lambda(\sqrt{n})$ amortized computation and communication cost using FHE. Zhou et al.~\cite{pir-ec23} and Lazzaretti and Papamanthou~\cite{lazza-pir} further improve the communication cost to $\widetilde{O}_\lambda(1)$, relying on privately puncturale PRF. 
%We introduce a single-server scheme that does not rely on the heavy cryptographic primitives, i.e., privately puncturable PRF in~\Cref{sec:polylog}.

All these aforementioned schemes require cryptographic assumptions that imply PKC, such as DDH and LWE.
Recnetly, Piano~\cite{piano} and a subsequent work from Mughees, I, and Ren~\cite{lingrenpir} introduce concretely efficient schemes based on the OWF assumption, while achieving optimal client-storage and online communication tradeoff (both are $\widetilde{O}_\lambda(\sqrt{n})$) and $\widetilde{O}_\lambda(\sqrt{n})$ amortized communication cost per query.
Our work preserves the optimal tradeoff and improves the online communication cost to $\widetilde{O}_\lambda(n^{1/4})$.


Doubly Efficient PIR (DEPIR), also known as the ``server-proprocessing model'', requires the server
first preprocesses the database and store additional bits, and subsequently, the server may answer queries with sublinear computation time. 
Earlier DEPIR schemes~\cite{depir00,depir01} rely on non-standard assumptions such as oblivious locally decodable code (OLDC) or virtual-blackbox obfuscation (VBB). 
Recently, a breakthrough work from Lin, Mook and Wichs~\cite{weikai-depir} constructs a DEPIR scheme based on the standard Ring-LWE assumption. 
